% The next line makes it so it compiles the main file if I hit the compile button
% when editing this file in TeXworks.
% !TEX root = ../ActiveIntroToDiscreteMathAndAlgorithms.tex

\chapter{Motivation}\label{ch:motivation}

The purpose of a discrete mathematics course in the computer science curriculum is
to give students a foundation in some of the mathematical concepts that are
foundational to computer science.   By ``foundational,'' we mean both
that the field of computer science was built upon (some of) them and 
that they are used to varying degrees in the study of the more advanced 
topics in computer science.

Computer science students sometimes complain about taking a discrete mathematics course.
They do not understand the relevance of the material to the rest of the
computer science curriculum or to their future career.
This can lead to lack of motivation.
They also perceive the material to be difficult.  
%Let's discuss these problems in a little more detail.

To be honest, some of the topics are difficult.  But the majority of the
material is very accessible to most students.  
One problem is that learning discrete mathematics takes effort, 
and when something doesn't sink in instantly, some students give up too quickly.  
The perceived difficulty together with a lack of motivation
can lead to lack of effort, which almost always leads to failure.
Even when students expend effort to learn, they can let their perceptions get the
best of them.  If someone believes something is hard or that they can't do it, it
often leads to self-fulfilling prophecy.  This is perhaps human nature.  
On the other hand, if someone believes that they can learn the material and solve the
problems, chances are they will.  The bottom line is that a positive attitude can
go a long way.  

%Success in learning discrete mathematics depends heavily on motivation.

This book was written in order
to ensure that the student {\em has to} expend effort while reading it.  The idea is 
that if you are allowed to simply read but not required to interact with the material,
you can easily read a chapter and get nothing out. 
For instance, your brain can go on `autopilot' when something doesn't sink in 
and you might get nothing out of the remainder of your time reading.  
By requiring you to solve problems and answer questions as you read, 
your brain is forced to stay engaged with the material. 
In addition, when you incorrectly solve a problem, you know immediately, 
giving you a chance to figure out 
what the mistake was and correct it before moving on to the next topic.
When you correctly solve a problem, your confidence increases.  
We strongly believe that this
feature will go a long way to help you more quickly and thoroughly learn the 
material, assuming you use the book as instructed.

What about the problem of relevance?  In other words, what is the connection
between discrete mathematics and other computer science topics?
There are several reasons that this connection is unclear to students.  
First, we don't always do a very good job of making the connection clear.  
We teach a certain set of topics because it is the
set of topics that has always been taught in such a course.   
We don't always think about the connection ourselves, 
and it is easy to forget that this connection is incredibly important to students.  
Without it, students can suffer from a lack of motivation to learn the material.

The second reason the connection is unclear is because one of the goals of such a course 
is simply to help students to be able to think mathematically.  As they continue in their
education and career, they will most certainly use some of the concepts they learn, yet they
may be totally unaware of the fact that some of their thoughts and ideas are based on 
what they learned in a discrete mathematics course.  Thus, although the students gain a 
benefit from the course, it is essentially impossible to convince them of this during
the course.

The third reason that the connection is unclear is that given the time constraints, 
it is impossible to provide all of the foundational mathematics that is relevant
to the advanced computer science courses {\em and} make the connection to those advanced topics
clear.  Making these connections would require an in-depth discussions of the advanced 
topics.  
%Do you see the problem?  
The connections are generally made, either implicitly or
explicitly, in the courses in which the material is needed.

This book attempts to address this problem by making connections to one set 
of advanced topics--the design and analysis of algorithms.  This is
an ideal application of the discrete mathematics topics since many of them 
are used in the design and analysis of algorithms.  We also do not have to go out of our
way too far to provide the necessary background, as we would if we attempted to make connections
to topics such as networking, operating systems, architecture, artificial intelligence, database,
or any number of other advanced topics.  As already mentioned, 
the necessary connections to those topics will be made
when you take courses that focus on those topics.

The goal of the rest of this chapter is to further motivate you to want to learn the topics
that will be presented in this book.  We hope that after reading it you {\em will} be more
motivated.  For some students, the topics are interesting enough on their own, whether or not they
can be applied elsewhere.  For others, this is not the case.  
One way or another, you must find motivation to learn this material.


\section{Some Problems}

In this section we present a number of problems for you to attempt to solve.
You should make an honest attempt to solve each.  We suspect that most readers will be
able to solve at most a few of them, and even then will probably not use the most
straightforward techniques.
On the other hand, after you have finished this book 
you should be able to solve most, if not all of them, with little difficulty.

There are two main reasons we present these problems to you now.  
First, we hope they help you gauge your learning.
That is, we hope that you {\em do} experience difficulty trying to to solve them now,
but that when you revisit them later, they will seem much easier.
Second, we hope they provide some motivation for you to learn the content.  
Although all of these problems may not interest you, 
we hope that you are intrigued by at least some of them.


\begin{description}
\item[Problem A:]
The following algorithm supposedly computes the sum of the first $n$ integers.  
Does it work properly?  If it does not work, explain the problem and fix it.
\begin{code}
    sum1ToN(int n) {
        return n + sum1ToN(n-1);
    }
\end{code}



\item[Problem B:]
The Mega Millions lottery involves picking five different numbers from 1 to 56, 
and one number from 1 to 46.  
I purchased a ticket last week and was surprised when none of my six numbers matched.  
Should I have been surprised?  
What are the chances that a randomly selected ticket will match none of the numbers?


\item[Problem C:]
I programmed an algorithm recently to solve an interesting problem.  
The input is an array of size $n$.  
When $n=1$, it took 1 second to run.  When $n=2$, it took 7 seconds.  When $n=3$, it took 19 seconds.  
When $n=4$, it took 43 seconds.  
Assume this pattern continues.
\begin{lenum}
\item
How large of an array can I run the algorithm on in less than 24 hours?  
\item
How large can $n$ be if I can wait a year for the answer?
\end{lenum}



\item[Problem D:]
Is the following a reasonable implementation of the {\sc Quicksort} algorithms?
In other words, is it correct, and is it efficient?
(Notice that the only difference between this and the standard algorithm is that
this one is implemented on a \verb+LinkedList+ rather than an \verb+array+.)
\begin{code}
	Quicksort(LinkedList A,int l,int r) {
	   if(r > l) {
	      int p = RPartition(A,l,r);
	      Quicksort(A,l,p-1);
	      Quicksort(A,p+1,r);
	   }
	}
	
	int RPartition(LinkedList A,int l,int r) {
	    int piv=l+(rand()%(r-l+1));
	    swap(A,l,piv);
	    int i = l+1;
	    int j = r;
	    while (1) {
	        while (A.get(i) <= A.get(l) && i<r) 
	            i++;
	        while (A.get(j) >= A.get(l) && j>l) 
	            j--;
	        if (i >= j) {
	            swap(A,j,l);
	            return j;
	        } else {
	            swap(A,i,j);
	        }
	    }
	}

	void swap(LinkedList A, index i, index j) {
	    int temp = A.get(i);
	    A.set(i,A.get(j));
	    A.set(j,temp);    
	}
\end{code}



\item[Problem E:]
I have an algorithm that takes two inputs, $n$ and $m$.  
The algorithm treats $n$ differently when it is less than zero, 
between zero and 10, and greater than 10.  
It treats $m$ differently based on whether or not it is even.  
I want to write some test code to make sure the algorithm works properly for all possible inputs.  
What pairs $(n,m)$ should I test?  
Do these tests guarantee correctness?  Explain.



\item[Problem F:]
Can the following code be simplified?  If so, give equivalent code that is as simple as possible.
\begin{code}
    if ((!x.size() <=0  && x.get(0) != 11) || x.size() > 0) 
    {
        if(!(x.get(0)==11 && (x.size() > 13 || x.size() < 13)) 
            && (x.size() > 0 || x.size() == 13)) 
        {
            //do something
        }
    }
\end{code}

\item[Problem G:]
In how many ways may we write the number $19$ as the sum of three
positive integer summands? Here order counts, so, for example, $1 +
17 + 1$ is to be regarded different from $17 + 1 + 1$.



\item[Problem H:]
Consider the \verb+stoogeSort+ algorithm given here:
\begin{code}
  void stoogeSort(int[] A,int L,int R){
      if(R<=L) { // Array has at most one element so it is sorted
          return;
      }
      if(A[R]<A[L]) {       
          int temp = A[L];  // Swap first and last element
          A[L] = A[R];      // if they are out of order
          A[R] = temp;
      }
      if(R-L>1){ // If the list has at least 3 elements
          int third=(R-L+1)/3;
          stoogeSort(A,L,R-third); // Sort first two-thirds
          stoogeSort(A,L+third,R); // Sort last two-thirds
          stoogeSort(A,L,R-third); // Sort first two-thirds again
      }
  }
\end{code}
\begin{lenum}
\item
Does \verb+stoogeSort+ correctly sort an array of integers?
\item
Is \verb+stoogeSort+ a good sorting algorithm?  Specifically, how long does it
take, and how does it compare to other sorting algorithms?
\end{lenum}






\item[Problem I:]
A cryptosystem was recently proposed.  
One of the parameters of the cryptosystem is a nonnegative integer $n$, 
the meaning of which is unimportant here.  
What is important is that someone has proven that the system is insecure for 
a given $n$ if there is more than one integer $m$ such that $2\cdot m \leq n \leq 2\cdot (m+1)$. 
\begin{lenum}
\item 
For what value(s) of $n$, if any, can you prove or disprove that  
there is more than one integer $m$ such that $2\cdot m \leq n \leq 2\cdot (m+1)$?
\item
Given your answer to (a), does this prove that the cryptosystem is either secure or insecure?
Explain.
\end{lenum}



\item[Problem J:]
A certain algorithm takes a positive integer, $n$, as input.
The first thing the algorithm does is set $n=n\bmod 5$. 
It then uses the value of $n$ to do further computations.
One friend claims that you can fully test the algorithm using just the
inputs $1$, $2$, $3$, $4$, and $5$.  Another friend claims that the
inputs $29$, $17$, $38$, $55$, and $6$ will work just as well.
A third friend responds with ``then why not just use $50$, $55$, $60$, $65$, and $70$?
Those should work just as well as your stupid lists.''
A fourth friend claims that you need many more test cases to be certain.
A fifth friend says that you can never be certain no matter how many test
cases you use.  Which friend or friends is correct?  Explain.


\item[Problem K:]
Write an algorithm to swap two integers without using any extra storage.
(That is, you can't use any temporary variables.)



\item[Problem L:]
You are at a party with some friends and one of them claims 
``I just did a quick count, and it turns out that at this party,
there are an odd number of people who have shaken hands with an odd number of other people.''
Can you prove or disprove that this friend is correct?



\item[Problem M:]
Recall the {\em Fibonacci sequence}, defined by the recurrence relation
$$f_n=\left\{
\begin{array}{ll}
0&\mbox{ if $n$=0}\\
1&\mbox{ if $n$=1}\\
f_{n-1}+f_{n-2}&\mbox{ if }n>1.\\
\end{array}
\right.
$$
So $f_2=1$, $f_3=2$, $f_4=3$, $f_5=5$, $f_6=8$, etc.
\begin{lenum}
\item
One friend claims that the following algorithm is an
elegant and efficient way to compute $f_n$.
\begin{code}
int Fibonacci(int n) {
    if(n <= 1) {
        return(n);
    } else { 
        return(Fibonacci(n-1)+Fibonacci(n-2));
    }
}
\end{code}
Is he right?  Explain.
\item
Another friend claims that he has an algorithm that computes $f_n$ that
takes constant time--that is, no matter how large $n$ is, it always takes the
same amount of time to computer $f_n$.  Is it possible that he has such an algorithm?
Explain.
\end{lenum}




\item[Problem N:]
You need to settle an argument between your boss (who can fire you) and your professor (who can fail you).  
They are trying to decide who to invite to the Young Accountants Volleyball League.  
They want to invite freshmen who are studying accounting and are over 6 feet tall.  
They have a list of everyone they could potentially invite.
\begin{enumerate}
\item
Your boss says they should make a list of all freshmen, 
a list of all accounting majors, and a list of everyone over 6 feet tall.  
They should then combine the lists (removing duplicates) and invite those on the combined list. 
\item
Your professor says they should make a list of everyone who is not a freshman, 
a list of anyone who does not do accounting, and a list of everyone who is 6 feet tall or less. 
They should make a fourth list that contains everyone who is on all three of the prior lists.  
Finally, they should remove from the original list everyone on this fourth list, 
and invite the remaining students.
\end{enumerate}
Who is correct? Explain.


\end{description}
